\section{Data and descriptive statistics}
 \label{sec:data}

\noindent
Our primary data source is the Moody's Industrial Manuals, from which we hand-collect annual balance sheet and income statement items for the period 1925--1940. Because some of our variables are growth rates, our sample period spans 1926--1940. The sample includes industrial public firms with coverage in the Center for Research in Securities Prices (CRSP) database, yielding a final sample of 7,074 firm-year observations with 558 unique firms.\footnote{We restrict our sample to publicly listed firms since equity market value is necessary to study investment under the lens of $Q$-theory, as discussed in Section \ref{sec:concept}.} Because our data are drawn from the Moody's Industrial Manuals, regulated industries (transportation, communication, and utilities) and financial firms---which are covered in separate Moody's manuals---are not part of our sample.

Our main variable of interest is net investment, defined as the annual growth rate in fixed capital stock, measured using the book value of property, plant, and equipment reported on the balance sheet. Our main explanatory variable, $d_{j,t}$, measures firm $j$'s exposure to the gold clause at time $t$ as the fraction of its long-term liabilities containing the clause:
\begin{equation}
\label{eq:d}
d_{j,t} = \frac{\text{\normalsize{Gold-denominated long-term liabilities of firm $j$ in year $t$}}}{\text{\normalsize{Long-term liabilities of firm $j$ in year $t$}}}.
\end{equation}
We define the long-term liabilities of a firm as the sum of the par value of these three components: corporate bonds, preferred shares, and bank loans. To construct $d_{j,t}$, we collect individual security-level information on the presence of the gold clause.\footnote{We collected bond-by-bond gold clause information for the period 1930--1935; our frozen measure $\tilde{d}_{j,t}$, introduced below, directly reflects these security-level data. Since the overwhelming majority of corporate bonds contained the gold clause, and given the cost of security-level data collection, we approximate gold-denominated liabilities using the book value of corporate bonds for years before 1930.} As noted in the previous section, the gold clause was included in most corporate bonds (around 97\% of them) and absent from all preferred shares. For bank debt, the manuals do not report exposure to the gold clause. However, at the time of the abrogation of the gold clause bank debt accounts for less than 2\% of long-term liabilities in the pre-abrogation period (see Table \ref{tab:sum_stats_d}). In the rest of the paper, we assume that bank debt did not include the clause. Our results are not sensitive to this assumption. Finally, firms without long-term liabilities are assigned $d_{j,t} = 0$.

Besides the main variables described above, we also collect balance sheet variables such as total assets, book equity, current assets, cash and marketable securities, along with net income from the income statement. These variables serve as control variables in some of our empirical specifications and allow us to construct standard financial ratios used in the investment literature. Appendix \ref{secapp:vars} contains detailed information on the construction of our variables.

Table \ref{tab:sum_stats_d} reports summary statistics separately for firms with no gold clause exposure ($d_{j,t} = 0$) and those with positive exposure ($d_{j,t} > 0$) across three key periods: 1926--1932 (Panel A), 1933--1934 (Panel B), and 1935--1940 (Panel C). In the pre-abrogation period, the average $d_{j,t}$ among exposed firms is 62\%.\footnote{The measure $d_{j,t}$ can slightly exceed the ratio of corporate bonds to long-term liabilities because $d_{j,t}$ is calculated using the reported amount outstanding of each bond, whereas Corp.\ bonds/LTL is based on balance sheet information. For most firms the two values coincide, but minor discrepancies between the book value of corporate bonds and the reported amount outstanding can arise.} Exposed firms are significantly larger (log assets of 17.42 vs.\ 16.77), more levered (book leverage of 42\% vs.\ 26\%; market leverage of 48\% vs.\ 28\%), hold less cash (12\% vs.\ 18\% of assets), generate lower net income (4\% vs.\ 6\% of assets), rely less on bank debt (1\% vs.\ 2\% of LTL), and pay out less to equity holders ($-$2\% vs.\ 13\% of common stock).

The descriptive statistics provide preliminary support for the view that firms with gold exposure cut investment by more during the period of legal uncertainty. From the pre-abrogation period to the litigation period, the average investment rate of exposed firms falls by 11 percentage points (from 5\% to $-$6\%), compared with a 6 percentage point decline for unexposed firms (from 3\% to $-$3\%). However, these descriptive patterns should be interpreted with caution, as pre-existing differences between exposed and unexposed firms may also drive investment decisions.\footnote{Internet Appendix Tables \ref{tabapp:summary_d_1} and \ref{tabapp:summary_d_0} report detailed summary statistics separately for $d_{j,t} > 0$ and $d_{j,t} = 0$ firms.} Our formal empirical tests will control for these differences, with Section \ref{sec:cont} further examining the sensitivity of our results to various control specifications.

Beyond observable differences, a separate challenge is that gold exposure is not randomly assigned across firms. In particular, Britain's departure from the gold standard in 1931 may have led firms to anticipate similar policy changes in the United States, prompting them to strategically adjust their gold clause exposure. For example, if only financially unconstrained firms were able to reduce their gold-denominated liabilities prior to the abrogation, the remaining firms with high gold exposure would disproportionately consist of constrained firms, potentially biasing the estimated effect of gold exposure on investment by conflating it with the impact of pre-existing financial constraints.

We investigate this concern by constructing a balanced panel of 157 firms that had corporate bonds outstanding at the end of 1930---the year before Britain left the gold standard---and track their outstanding bonds until the end of 1934. Specifically, we examine whether firms adjusted the number of outstanding gold-denominated bonds or reduced the share of gold clause liabilities in their capital structure prior to the 1933 abrogation. Table \ref{tab:bond_stats} reports our findings. On the extensive margin, firms made minimal adjustments: of the 157 firms with corporate bonds outstanding, 150 had gold clause exposure at the end of 1930, and 144 of these still had gold-denominated bonds by the end of 1932. Similarly, the total number of gold-denominated bonds outstanding declined only modestly from 250 to 240 over the same period. On the intensive margin, we find no evidence of strategic reduction in exposure levels. Among firms with positive gold clause exposure, the average value of $d_{j,t}$ remained virtually unchanged at 0.57 in 1930 and 0.58 in 1932. These patterns suggest that firms either failed to anticipate the eventual abandonment of the gold standard or faced prohibitive costs to adjust their exposure.

As discussed in Section \ref{sec:historical_background}, both President Hoover and presidential candidate Roosevelt defended the standard throughout 1932, with Roosevelt explicitly denying he would alter the value of gold if elected. This apparent bipartisan commitment may have convinced firm managers that abandoning gold had little to no political support. Moreover, even if some managers perceived the risk, market conditions arguably made debt restructuring prohibitively expensive. As Figure \ref{fig:cfc_bond_quotes} in the Internet Appendix shows, industrial bonds traded at steep discounts to par value in 1931, which implies that most call options were deeply out-of-the-money. Open-market repurchases would have been similarly costly, as buying large quantities of bonds would likely have moved prices against the firms.

Nonetheless, to address potential endogeneity from post-1931 adjustments in gold exposure, we construct an alternative measure, $\tilde{d}_{j,t}$, that freezes each firm's exposure at its 1930 level for all subsequent years. It is formally defined as:
\begin{equation}
\label{eq:tilde_d}
\tilde{d}_{j,t} = \begin{dcases}
d_{j,t}, & \text{if } t \leq 1930 \\
d_{j,1930}, & \text{if } t > 1930
\end{dcases}.
\end{equation}
By construction, $\tilde{d}_{j,t}$ captures a firm's gold exposure before Britain's policy shift could have prompted strategic adjustments. The evidence presented above suggests this approach is conservative: firms made minimal changes to their gold-denominated liabilities even after 1931, so $\tilde{d}_{j,t}$ and $d_{j,t}$ are highly correlated (as shown in the last column of Table \ref{tab:bond_stats}). With these data and measures in hand, we now turn to our formal empirical analysis.\footnote{The Internet Appendix provides additional summary statistics based on $\tilde{d}_{j,t}$. Table \ref{tab:sum_stats_tilde_d} is the analog of Table \ref{tab:sum_stats_d} using $\tilde{d}_{j,t}$ as the grouping variable instead of $d_{j,t}$. Tables \ref{tabapp:summary_I_0} and \ref{tabapp:summary_I_1} report a broader set of summary statistics separately for $\tilde{d}_{j,t} > 0$ and $\tilde{d}_{j,t} = 0$ firms. Tables \ref{tabapp:summary_I_smalld} and \ref{tabapp:summary_I_larged} further split the $\tilde{d}_{j,t} > 0$ sample into high- and low-$\tilde{d}_{j,t}$ firms and show that high-$\tilde{d}_{j,t}$ firms are slightly smaller and less levered than low-$\tilde{d}_{j,t}$ firms. Table \ref{tabapp:correlation} reports correlations between $\tilde{d}_{j,t}$ and firm characteristics.}